\section{Fungsi}
\subsection{Pengenalan Fungsi}
\begin{definition}[Fungsi]
  Fungsi adalah aturan yang mengubah nilai variabel menjadi nilai angka, biasanya ditulis dengan notasi $f(x)$ dibaca ``$f$ pada $x$'', artinya aturan fungsi $f$ berlaku untuk variabel $x$.
\end{definition}

\begin{example}[Fungsi]
  \hfill \\
  Misalkan fungsi:
  \[ f(x) = x^2 + 2x - 7 \]
  \begin{enumerate}
  \item Evaluasi $f(1)$:
    \[
    \begin{aligned}
      f(1) &= (1)^2 + 2(1) - 7 \\
      &= 1 + 2 - 7 \\
      &= -4
    \end{aligned}
    \]
  \item Evaluasi $f(2)$:
    \[
    \begin{aligned}
      f(2) &= (2)^2 + 2(2) - 7 \\
      &= 4 + 4 - 7 \\
      &= 1
    \end{aligned}
    \]
  \item Evaluasi $f(-2)$:
    \[
    \begin{aligned}
      f(-2) &= (-2)^2 + 2(-2) - 7 \\
      &= 4 - 4 - 7 \\
      &= -7
    \end{aligned}
    \]
  \end{enumerate}
\end{example}

\subsection{Penjumlahan Fungsi Aljabar}
\begin{definition}[Penjumlahan Fungsi]
  Jika diberikan 2 buah fungsi $f(x)$ dan $g(x)$, maka penjumlahan fungsi didefinisikan dengan rumus:
  \[ (f + g)(x) = f(x) + g(x) \]
  Artinya, jumlahkan bentuk aljabar dari kedua fungsi dengan cara menjumlahkan suku yang sejenis (misalkan $x^2$ dengan $x^2$, $x$ dengan $x$, konstanta dengan konstanta, dengan memperhatikan tanda positif dan negatif).
\end{definition}

\begin{example}[Penjumlahan Fungsi]
  \hfill
  \begin{enumerate}
  \item Jika $f(x) = 4x + 6$ dan $g(x) = 2x - 3$:
    \[
    \begin{aligned}
      (f + g)(x) &= (4x + 6) + (2x - 3) \\
      &= 4x + 2x + 6 - 3 \\
      &= 6x + 3
    \end{aligned}
    \]
  \item Jika $f(x) = 3x^2 - x - 4$ dan $g(x) = -x^2 + 5x - 3$:
    \[
    \begin{aligned}
      (f + g)(x) &= (3x^2 - x - 4) + (-x^2 + 5x - 3) \\
      &= 3x^2 - x^2 - x + 5x - 4 - 3 \\
      &= 2x^2 + 4x - 7
    \end{aligned}
    \]
  \end{enumerate}
\end{example}
